\section{Вступ}
Соціальні конфлікти дедалі частіше залишають паралельні записи в каналах, гілках коментарів, новинних повідомленнях, заявах установ і коротких аудіовізуальних фрагментах. Ці записи слугують різним аналітичним завданням: зафіксувати, що сказало джерело, простежити поширення твердження, виявити повторений текст, упорядкувати можливі події та оцінити, чи підтримують наявні докази певне звинувачення. Система, яка змішує ці завдання, може перетворити правдоподібну оповідь на нібито перевірений опис. Скопійоване твердження не є новим свідком, час публікації не є часом запису, а статистичне відхилення від моделі часу не вказує на організатора.

Кризова інформатика давно наголошує на складності отримання корисної інформації з неоднорідних і зашумлених потоків соціальних мереж \citep{imran2015crisis}. Обчислювальна соціальна наука додає питання доступу до платформ, вибірки, вимірювання та етики \citep{Olteanu2019,VicentePaula2023}. Цифрове розслідування за відкритими джерелами потребує збереження й контекстуальної перевірки матеріалів, зокрема явного опису невизначеності \citep{berkeleyProtocol2022}. Ці вимоги стають особливо важливими, коли програмне забезпечення поєднує традиційний аналіз із генеративним ШІ. Вправний синтез може приховувати невизначеність, а не усувати її; навіть генерація з цитуванням потребує окремої оцінки того, чи підтримують цитати сам текст \citep{gao2023alce}.

Social розв'язує цю інтеграційну проблему як дослідницький програмний фреймворк. Його практична мета --- допомогти аналітику перейти від зібраних цифрових слідів до придатних для перегляду проміжних записів та обмежених аналітичних результатів. Реалізація містить скрипти отримання даних для окремих платформ, аналіз інформаційних хвиль, зв'язаний реєстр доказів, правила пошуку й перевірки, обробку типізованих подій, локальні медіаконвеєри та валідатори звітів. Ці компоненти суттєво різняться за зрілістю, інтегрованістю й оціненістю. Опис їх як єдиного рівномірно автономного дослідника приховав би ці відмінності.

Стаття ставить питання: \emph{як модульна, доказово обґрунтована програмна система може інтегрувати неоднорідні цифрові сліди для автоматизованого аналізу соціальних конфліктів та інформаційних інцидентів, зберігаючи походження даних, невизначеність і відтворюваність?} Ми вивчаємо наявну реалізацію, а не пропонуємо новий алгоритм. Аналіз фіксує ревізію Git, досліджує шляхи виконання, звіряє числові артефакти, повторно запускає обмежені обчислення й навмисно перевіряє випадки семантичних збоїв. Кейс Петропавлівської Борщагівки дає реалістичні робочі навантаження, але не є незалежною перевіркою узагальнюваності. Сам інцидент вивчено в окремому рукописі проєкту \citep{pylypenko2026case}; ця стаття розглядає програмне забезпечення та межі його доказовості.

\subsection{Внесок роботи}
Внесок роботи складається з чотирьох частин. По-перше, подано інтеграційну карту, засновану на коді, яка розмежовує повторно використовувані процесори, збирачі платформ, експериментальні евристики, кураторські збирачі кейсу та скрипти формування звітів. По-друге, описано реалізовану модель доказів і розрізнення технічного походження, лінії походження інформації, незалежності джерел та семантичної обґрунтованості. По-третє, наведено перевірки відтворюваності для виявлення схожості, умовної рандомізації, структурних контрактів та обмеженого офлайн-процесу. По-четверте, показано конкретні семантичні контрприклади, виміряно на них оприлюднену ревізію правила й визначено межі продемонстрованої автономності системи. Цей внесок сприяє відповідальнішому використанню й оцінюванню аналітичного програмного забезпечення, але не встановлює нового рівня техніки у перевірці фактів, виділенні подій чи атрибуції.
